% ECONOMICS_PROBLEM: {"id": "H55-336", "title": "Public pensions and longevity risk", "jel_code": "H55", "source_id": "OP-0479"}
% !TeX program = xelatex
\documentclass[11pt,a4paper]{article}
\usepackage[margin=23mm]{geometry}
\usepackage{fontspec}
\setmainfont{Latin Modern Roman}
\usepackage{amsmath,amssymb,mathtools}
\usepackage{xcolor,enumitem,booktabs,longtable,array,xurl,hyperref}
\definecolor{muted}{HTML}{53636C}
\hypersetup{colorlinks=true,urlcolor=blue,linkcolor=blue}
\setlength{\parindent}{0pt}
\setlength{\parskip}{5pt}
\newcommand{\R}{\mathbb R}
\newcommand{\E}{\mathbb E}
\newcommand{\N}{\mathbb N}
\newcommand{\ind}{\mathbf 1}
\newcommand{\Var}{\operatorname{Var}}
\newcommand{\Cov}{\operatorname{Cov}}
\newcommand{\supp}{\operatorname{supp}}
\DeclareMathOperator*{\argmin}{arg\,min}
\DeclareMathOperator*{\argmax}{arg\,max}
\newcommand{\entrymeta}[3]{{\sffamily\small\color{muted}\textbf{\texttt{#1}}\quad|\quad #2\par #3\par}}
\newcommand{\Problemref}[1]{\texttt{#1}}
\newcommand{\JELref}[2]{\texttt{#2} (JEL \texttt{#1})}
\begin{document}
\section*{Public pensions and longevity risk}
\textbf{Problem ID:} \texttt{H55-336}\quad \textbf{JEL:} \texttt{H55}\par
% BEGIN ECONOMICS STATEMENT
\label{problem:OP-0479}
{\sffamily\small\color{muted}Original subject group: Taxation and social insurance\par}
\label{definitions:problem:30007761}
\textbf{Audit classification:} Background; proposed extension. The published author-hosted four-page PDF verifies authors, 2026 volume/DOI and the solved benchmark. Repaired undiscounted death bequests, survival indexing and time-model and wealth qualifications. The source does not list this heterogeneous extension as open.
\subsection*{Definitions and precise research target}
Fix a retirement cohort of types \(\theta\sim F\), observable immutable category \(x_\theta\), initial wealth \(w_\theta\geq0\), utility \(u_\theta\), bequest utility \(v_\theta\), discount \(0<\beta_\theta<1\), and death time \(\tau\in\{1,2,\ldots\}\cup\{\infty\}\). The survival function is \(q_\theta(t)=\Pr_\theta(\tau>t)\); death risk and its possible dependence on policy are explicitly specified. In the baseline take survival as exogenous, no private annuities, no borrowing, and deterministic savings return \(r_a>-1\). A pension \(b_t(x)\geq0\) is payable only while alive. Conditional on survival, wealth satisfies
\[
 c_t+\frac{a_{t+1}}{1+r_a}=a_t+b_t(x_\theta),\qquad
 a_0=w_\theta,\quad a_t,c_t\geq0.
\]
If death occurs before date \(t\)'s consumption, the estate is \(a_t\); any accidental bequest enters \(v_\theta\). Specify a no-Ponzi/transversality condition and assume finite expected utility. Let \(J_\theta(b)\) be the supremum over feasible contingent consumption/savings plans of
\[
 \mathbb E_\theta\!\left[\sum_{t<\tau}\beta_\theta^tu_\theta(c_t)
 +\mathbf1_{\{\tau<\infty\}}\beta_\theta^\tau v_\theta(a_\tau)\right].
\]
With fixed social weights \(\omega_\theta\), government discount rate \(r_g>-1\), budget \(R\), and admissible pension class \(\mathcal B\), define
\[
 \sup_{b\in\mathcal B}\int\omega_\theta J_\theta(b)\,dF(\theta),\qquad
 \int\sum_{t\geq0}\frac{q_\theta(t)b_t(x_\theta)}{(1+r_g)^t}\,dF(\theta)\leq R.
\]
State feasibility, convergence and attainment separately. Characterize which heterogeneity and bequest restrictions preserve a deferral pattern, and derive identifying requirements for welfare comparisons. Manipulable eligibility categories require additional incentive constraints. Private annuity access needs a contract/menu/price model, not only the phrase `private market'. The cited source solves its representative-retiree benchmark: its representative benchmark has different discrete and continuous-time qualifications. This heterogeneous extension is editorial and is not certified open by the benchmark theorem.
\subsection*{Variants and assumption changes}
\begin{enumerate}
\item \textbf{Representative versus heterogeneous retiree (editorial).} With a common type and no bequests, use the published benchmark; with type-dependent survival and wealth and common observable categories, optimize the cohort objective. Redistribution and longevity insurance interact.
\item \textbf{No annuities versus private contracts (editorial).} Add a declared annuity menu, timing, solvency constraints and prices; if prices are endogenous, define selection and insurer zero-profit conditions. Public benefits can crowd out private coverage.
\item \textbf{Discrete finite versus continuous time (source benchmark distinction).} In the paper's two-period benchmark, with initial wealth \(w\), budget \(S\) and survival \(p>0\), terminal-only benefits require \(w\geq S/p\). If \(w<S/p\), first-period benefits are \((S-pw)/(1+p)>0\). Longer finite horizons with binding borrowing constraints can have a deferred flat stream and one smaller top-up payment. The continuous-time infinite-horizon benchmark has a deferred level stream. These shape results assume the source's no-interest, no-bequest representative setting.
\end{enumerate}
\subsection*{References and source audit}
Einav--Finkelstein (2026), author-hosted published PDF, Journal of Public Economics 256, 105601: sections 2--3 and conclusion verified. The representative benchmark is solved. A heterogeneous/private-market extension is proposed here, not asserted source-authored.
\begin{enumerate}\raggedright
\item\hypertarget{definitions:source:30007761:1}{} Liran Einav; Amy Finkelstein. 2026. \emph{On the Optimality of Deferred Public Annuities}. Published article, sections2–3 and conclusion; Journal of Public Economics256(2026),105601.
\url{https://web.stanford.edu/~leinav/pubs/JPubEc2026.pdf}.
DOI: \url{https://doi.org/10.1016/j.jpubeco.2026.105601}.
\end{enumerate}

\subsection*{Common conventions: Source mathematical conventions}

These conventions apply to each entry unless explicitly replaced.

\textbf{Models and identification.} A primitive model \(m\) specifies all probability laws, feasible choices, information, equilibrium and measurement rules used by its target. The admissible class \(\mathcal M\) is fixed independently of outcomes. If \(P_O(m)\) is its observable probability law and \(\tau(m)\) the target, its identified set at observable law \(P\) is
\[
 \mathcal I_\tau(P)=\{\tau(m):m\in\mathcal M,\ P_O(m)=P\}.
\]
Point identification means that this set is a singleton. An empty set signals incompatibility of data and assumptions. Coordinate bounds need not describe the joint set. Bounds are sharp only if every retained target is attainable or arbitrarily approachable by compatible models. Finite bounds require explicit restrictions on unobserved quantities; unrestricted confounding, hidden wealth, extrapolation or arbitrary equilibrium selection can yield uninformative sets.

\textbf{Causal interventions.} A potential outcome \(Y_i(p)\) is the outcome under a complete policy regime \(p\), including timing, financing, information and market spillovers. The target population and units are fixed before treatment. With interference, use the complete assignment vector or a justified exposure mapping; individual no-interference notation alone cannot identify an economy-wide intervention. A difference of mean potential outcomes does not require the cross-world joint distribution. Conditioning on post-treatment migration, survival, enrollment or employment changes the estimand and needs separate justification.

\textbf{Statistical statements.} An estimator, confidence set, forecast or test specifies a sampling law, model class and loss/coverage criterion. Uniform \((1-\alpha)\)-coverage means \(\inf_{m\in\mathcal M}\Pr_m\{\tau(m)\in C_n\}\geq1-\alpha\), or an explicitly asymptotic version. The whole parameter space trivially has coverage, so informativeness requires a stated width, risk or power objective. A forecasting improvement is relative to a named benchmark, information set and evaluation population.

\textbf{Equilibrium and optimization.} An equilibrium comprises feasible plans, optimality, beliefs where needed, budget constraints and all market/resource clearing. The primitives or a stated benchmark define each correspondence \(\mathcal E(p;m)\). Nonemptiness is assumed or proved, never inferred just from compact policy sets. A selected-equilibrium target uses a declared selection rule; a robust target quantifies over all permitted equilibria. Use suprema or infima unless attainment is established. An \(\varepsilon\)-optimal policy is feasible and within \(\varepsilon\) of the finite optimum value. Report an empty feasible set separately.

\textbf{Welfare and accounting.} Normative utility, interpersonal normalization, social weights, numeraire, discounting and the use of public receipts are primitives. Behavioral utility need not coincide with the welfare criterion. Household welfare includes ownership income and rents once; adding profits again to compensating variation generally double-counts them. A monetary equivalent is defined by an explicit transfer and price convention, with monotonicity and range conditions. Average effects, mechanism contributions, transfers and welfare losses are different targets. Infinite sums require convergence and dynamic feasible plans require terminal or no-Ponzi conditions.

\textbf{Source versions.} A working-paper date, revision date and journal publication year are distinguished. A DOI for a journal version is not automatically the DOI of an earlier manuscript. Locators refer to the stated version and printed pages where specified. A metadata-only check does not verify a claim about a full-text conclusion. Every entry's source subsection narrows the provenance claim accordingly.
% END ECONOMICS STATEMENT
\end{document}
